\documentclass[a4paper,oneside,12pt]{extarticle}

\usepackage[utf8]{inputenc}
\usepackage[english]{babel}
\usepackage{amssymb,amsthm,mathtools,enumitem,geometry,hyperref,booktabs}
\usepackage{graphicx}
\usepackage{xcolor}
\usepackage{tikz}
\usetikzlibrary{calc,patterns,decorations.pathreplacing}
\usepackage{tikz-cd}
\usepackage{pgfplots}
\usepgfplotslibrary{fillbetween}
\pgfplotsset{compat=1.18}
\usepackage{float}

\geometry{lmargin=15mm,rmargin=15mm,tmargin=10mm,bmargin=10mm,includeheadfoot}
\pagestyle{plain}
\frenchspacing

\theoremstyle{plain}
\newtheorem{proposition}{Proposition}
\newtheorem{lemma}[proposition]{Lemma}

\theoremstyle{definition}
\newtheorem{definition}[proposition]{Definition}
\newtheorem{example}[proposition]{Example}

\theoremstyle{remark}
\newtheorem{remark}[proposition]{Remark}

\newcommand*{\E}{\mathbb{E}}
\newcommand*{\Var}{\mathrm{Var}}
\newcommand*{\mP}{\mathbb{P}}
\newcommand*{\mR}{\mathbb{R}}
\newcommand*{\mbF}{\mathcal{F}}
\newcommand*{\ve}{\varepsilon}




\begin{document}
\begin{titlepage}
\begin{center}

{\footnotesize\scshape\color{gray}
Working Notes in Quantitative Risk and Modeling
}

\vfill

{\LARGE\bfseries Homoscedasticity and Its Testing:\\[4pt]
From Conditional Expectation to the Breusch--Pagan Test}

\vspace{36pt}

{\normalsize Boris Garbuzov}

\vfill

{\Large\bfseries Part~7~/~12}

\vspace{6pt}

{\Large\bfseries The Lagrange Multiplier Test:\\Why Score and LM Are One Thing}

\vfill

{\small July 2026}

\end{center}
\end{titlepage}

\setcounter{tocdepth}{2}
\tableofcontents
\newpage
\section{The Lagrange Multiplier Method and the LM Test}
\label{sec:lm}

\subsection{Basic Lagrange multipliers: why the score test carries the name LM}
\label{subsec:basic-lm}

We follow the Probabilist's notes of June~8, 2026.  Before connecting the score
test to the LM name, we recall the Lagrange multiplier method
itself --- purely as a calculus fact, no randomness, no testing yet.

\subsubsection*{Step 1: unconstrained maximisation}

Consider maximising $f(x,y)$ over the full domain $D\subseteq\mathbb{R}^2$.

\begin{figure}[H]
\centering
\includegraphics[width=0.60\textwidth]{fig_lm1.png}
\caption{Figure~1 (the Probabilist, June 8 2026).  Unconstrained problem:
$f(x,y)\to\max$ over $D$.  The red surface is $f$; the argmax lies
at the peak.}
\label{fig:lm1}
\end{figure}

\subsubsection*{Step 2: general subset constraint}

Now restrict to a subset $\tilde{D}\subseteq D$.  This can only lower
the maximum:
\[
  \tilde{D}\subseteq D \;\Longrightarrow\; \max_{\tilde{D}} f \leq \max_D f.
\]

\begin{figure}[H]
\centering
\includegraphics[width=0.60\textwidth]{fig_lm2.png}
\caption{Figure~2 (the Probabilist).  Constrained problem: $f(x,y)\to\max$
over $\tilde{D}\subseteq D$ (shaded region).  The constrained maximum
is at most the unconstrained one.  This general constraint has no
special relation to LM; it is shown for contrast.}
\label{fig:lm2}
\end{figure}

\subsubsection*{Step 3: equality constraint via a smooth function}

Now suppose $\tilde{D}$ is the zero-level set of a smooth function $g$:
\[
  \tilde{D} = \{(x,y)\mid g(x,y)=0\}.
\]
This reduces the dimension of the feasible set by one.

\begin{figure}[H]
\centering
\includegraphics[width=0.60\textwidth]{fig_lm3.png}
\caption{Figure~3 (the Probabilist).  The green surface is $g(x,y)$; its
zero-level set $\tilde{D}=\{g=0\}$ is the curve on the floor.
The \emph{constraint qualification} $\nabla g\neq 0$ at the
solution is assumed throughout.}
\label{fig:lm3}
\end{figure}

\subsubsection*{Step 4: combining $f$ and the constraint}

Figure~4 overlays $f$ (red) and $g$ (green).  The constrained argmax
$B=(x^*,y^*,0)$ lies on the floor; the Lagrange multiplier value
$\lambda^*$ is the third coordinate of the saddle point
$A=(x^*,y^*,\lambda^*)$ of the Lagrangian.

\begin{figure}[H]
\centering
\includegraphics[width=0.60\textwidth]{fig_lm4.png}
\caption{Figure~4 (the Probabilist).  Combined picture:
$A = \max f\!\mid_{g=0} = (x^*,y^*,\lambda^*)$ and
$B = \operatorname{argmax} f\!\mid_{g=0} = (x^*,y^*,0)$.}
\label{fig:lm4}
\end{figure}

\subsubsection*{Step 5: the Lagrangian and the third dimension}

To find $(x^*,y^*)$ we do not enlarge $\tilde{D}$ within $\mathbb{R}^2$.
Instead we \emph{expand the ambient space} by adding a third
coordinate $\lambda$, going from $\mathbb{R}^2$ to $\mathbb{R}^3$.
On this extended space define the \emph{Lagrangian}:
\[
  \mathcal{L}(x,y,\lambda) = f(x,y) - \lambda\cdot g(x,y).
\]

\begin{figure}[H]
\centering
\includegraphics[width=0.65\textwidth]{fig_lm5.png}
\caption{Figure~5 (the Probabilist).  The orange hatching shows the
$\lambda$-``wall'' above the constraint curve.  The unconstrained
stationary point of $\mathcal{L}$ in $\mathbb{R}^3$ projects onto
the constrained argmax in $\mathbb{R}^2$.  By the constraint
qualification, $(x^*,y^*)$ is not a stationary point of $g$ itself.}
\label{fig:lm5}
\end{figure}

Setting all three partial derivatives of $\mathcal{L}$ to zero:
$\partial_\lambda\mathcal{L}=0$ recovers the constraint $g(x,y)=0$;
$\partial_x\mathcal{L}=0$ and $\partial_y\mathcal{L}=0$ give the
parallel-gradients condition $\nabla f = \lambda\,\nabla g$.

The value $\lambda^*$ at the solution is the \emph{Lagrange
multiplier}.  It measures the shadow price of the constraint: how
much the constrained maximum would increase if $g=0$ were relaxed to
$g=\varepsilon$.

\medskip\noindent
\textbf{Connection to the score test.}  In the likelihood context,
$f = \ell(\theta)$ and the constraint $g(\theta)=0$ encodes the null
hypothesis $H_0$.  The Lagrange multiplier $\lambda^*$ at the
restricted MLE $\hat\theta_0$ is proportional to the score
$\partial_\theta\ell(\hat\theta_0)$.  A large $|\lambda^*|$ means
the likelihood is straining against the constraint: the data push
away from $H_0$.  This is exactly what the score test measures,
which is why ``LM test'' and ``score test'' are two names for the
same statistic.

\subsection{Basic LM 2: gradient collinearity and level lines}
\label{subsec:lm_gradient}

We follow the Probabilist's notes of June 9, 2026.  This subsection continues
the Lagrange multiplier story with the geometric picture of gradient
collinearity, and connects it to the tangency of level lines.

\medskip
\noindent\textbf{Step 6: gradient collinearity condition.}
At the constrained optimum $(x^*, y^*)$, moving along the constraint
$g = 0$ cannot increase $f$.  This means $\nabla f$ must be
orthogonal to the constraint curve $\{g=0\}$.  But $\nabla g$ is
also orthogonal to $\{g=0\}$ (gradient is always perpendicular to
level sets).  Hence the two gradients must be parallel (collinear):
\[
  \nabla f(x^*, y^*) = \lambda^* \nabla g(x^*, y^*).
\]
To find $(x^*, y^*)$ we search along the green curve on the floor for
the point where $\nabla f \parallel \nabla g$.

\begin{figure}[H]
  \centering
  \includegraphics[width=0.72\textwidth]{fig_lm6.png}
  \caption{The Probabilist's sketch (9 June 2026): the 3-d picture with
  $\nabla f$ (black) and $\nabla g$ (green) marked at the constrained
  argmax $B=(x^*,y^*,0)$.  Point $A = (x^*,y^*,\lambda^*)$ is the
  saddle point of the Lagrangian.}
\end{figure}

\medskip
\noindent\textbf{Step 7: gradient collinearity in 2d.}
Projecting to the $(x,y)$ plane: the constraint $\{g=0\}$ is the
green ellipse.  At $(x^*,y^*)$ the vector $\nabla g$ is orthogonal
to the ellipse (outward normal).  The condition $\nabla f \parallel
\nabla g$ says $\nabla f$ points in the same direction.

\begin{figure}[H]
  \centering
  \includegraphics[width=0.60\textwidth]{fig_lm7.png}
  \caption{The Probabilist's sketch (9 June 2026): 2-d projection.  At
  $(x^*,y^*)$ the gradients $\nabla f$ and $\nabla g$ are collinear.
  $\nabla g$ is orthogonal to the level line of $g$ (green ellipse).}
\end{figure}

\medskip
\noindent\textbf{Step 8: raising the level line of $f$.}
We can cut the surface $f$ at the level $f(x^*,y^*)$ to obtain a
\emph{raised level line of $f$} — the set of points in 3-d at height
$f(x^*,y^*)$:
\[
  \{(x,y,f(x,y)) \mid f(x,y) = f(x^*,y^*)\}
  = \{(x,y,f(x^*,y^*)) \mid f(x,y) = f(x^*,y^*)\}.
\]

\begin{figure}[H]
  \centering
  \includegraphics[width=0.72\textwidth]{fig_lm8.png}
  \caption{The Probabilist's sketch (9 June 2026): the raised level line of $f$
  (red curve on the surface) passes through $A=(x^*,y^*,z^*)$ where $z^*=f(x^*,y^*)$.
  Its projection onto the floor gives the level line of $f$.}
\end{figure}

\medskip
\noindent\textbf{Step 9: projecting to the floor.}
Projecting the raised level line down to the $(x,y)$ plane gives the
ordinary level line of $f$:
\[
  \{(x,y) \mid f(x,y) = f(x^*,y^*)\}.
\]
This level line passes through $(x^*,y^*)$ by construction.

\begin{figure}[H]
  \centering
  \includegraphics[width=0.72\textwidth]{fig_lm9.png}
  \caption{The Probabilist's sketch (9 June 2026): the level line of $f$ (red,
  $f=f(x^*,y^*)$) shown on the floor together with the constraint
  curve $g=0$ (green).  Both pass through $(x^*,y^*)$.}
\end{figure}

\medskip
\noindent\textbf{Step 10: tangency of level lines.}
In 2-d, the level line of $f$ and the constraint curve $\{g=0\}$ are
\emph{tangent} at $(x^*,y^*)$.  Tangency is exactly the geometric
content of $\nabla f \parallel \nabla g$: both curves share the same
tangent direction at the contact point, because their normals
($\nabla f$ and $\nabla g$) coincide up to a scalar.

\begin{figure}[H]
  \centering
  \includegraphics[width=0.60\textwidth]{fig_lm10.png}
  \caption{The Probabilist's sketch (9 June 2026): level line of $f$
  ($f = f(x^*,y^*)$, red ellipse) and constraint curve $g=0$ (green
  ellipse) are tangent at $(x^*,y^*)$.  At the tangency point
  $\nabla f \parallel \nabla g$.}
\end{figure}

\noindent\textit{Remark.}
The Hessian matrix determines whether the constrained stationary
point is a maximum or minimum of $f$.  If there are $k$ equality
constraints $g_1=\cdots=g_k=0$, the Lagrangian acquires $k$
multipliers $\lambda_1,\ldots,\lambda_k$; the number of constraints
must not exceed the number of variables.  Inequality constraints lead
to the Kuhn--Tucker conditions, a generalisation not pursued here.

\subsection{Score as a gradient vector}

We follow the Probabilist's notes of June 10, 2026.  So far we have treated the
score $s(\theta,x)=\partial_\theta\ell(\theta,x)$ in the
one-dimensional case, where it is a scalar --- the slope of the
log-likelihood.  For a multivariate parameter $\theta\in\mathbb{R}^p$
the score becomes a \emph{gradient vector}:
\[
  S(\theta,x) = \nabla_\theta\,\ell(\theta,x)
  = \Bigl(\frac{\partial\ell}{\partial\theta_1},\,
           \frac{\partial\ell}{\partial\theta_2},\,\ldots,\,
           \frac{\partial\ell}{\partial\theta_p}\Bigr)^\top.
\]
The log-likelihood $\ell(\theta_1,\theta_2)$ is now a surface over the
parameter plane $(\theta_1,\theta_2)$, and the score $S$ is the
gradient field of that surface --- a vector field pointing in the
direction of steepest ascent at each $\theta$.

\begin{figure}[H]
  \centering
  \includegraphics[width=0.72\textwidth]{fig_sv1.png}
  \caption{The Probabilist's sketch (10 June 2026): the log-likelihood surface
  $\ell(\theta_1,\theta_2)$ (red) over the $(\theta_1,\theta_2)$
  plane.  The red lines are hatching of the surface, not gradient
  arrows.  The green arrow marks the gradient $\nabla\ell$ at a
  specific point $\theta$.}
\end{figure}

\noindent\textbf{The score as a map $\mathbb{R}^2\to\mathbb{R}^2$.}
For $p=2$, the score $S(\theta,x)$ is an operator that takes a point
$\theta=(\theta_1,\theta_2)$ in parameter space and returns the
two-dimensional gradient vector $(s_1,s_2)=(\partial_{\theta_1}\ell,
\partial_{\theta_2}\ell)$.  The left panel shows the domain
$(\theta_1,\theta_2)$ with a sample point; the right panel shows the
image $(s_1,s_2)$ in score space.

\begin{figure}[H]
  \centering
  \includegraphics[width=0.72\textwidth]{fig_sv2.png}
  \caption{The Probabilist's sketch (10 June 2026): the score as an operator
  $S\colon\mathbb{R}^2\to\mathbb{R}^2$.  Left: parameter space
  $(\theta_1,\theta_2)$.  Right: score space $(s_1,s_2)$.  At the
  MLE, $S=0$ (the score maps to the origin).}
\end{figure}

\noindent\textbf{The gradient field.}
Evaluating $S(\theta,x)$ at every $\theta$ in the parameter plane
produces a gradient vector field.  All arrows point toward the MLE
(the red dot), where the field vanishes: $S(\hat\theta,x)=0$.  The
level curves of $\ell$ (red ellipse) are orthogonal to the gradient
arrows at every point.

\begin{figure}[H]
  \centering
  \includegraphics[width=0.65\textwidth]{fig_sv3.png}
  \caption{The Probabilist's sketch (10 June 2026): gradient field of
  $\ell(\theta_1,\theta_2)$ over the parameter plane.  Green arrows
  show $\nabla\ell$ at a grid of points; all arrows converge on the
  MLE (red dot).  The red ellipse is a level curve of $\ell$.}
\end{figure}

\noindent\textbf{The two component surfaces $s_1$ and $s_2$.}
Each component $s_k(\theta,x)=\partial_{\theta_k}\ell(\theta,x)$ can
be drawn as a separate surface over $(\theta_1,\theta_2)$.  When
$\omega$ (a random draw from the sampling distribution) is plugged in
for $x$, the likelihood $\ell(\theta,\omega)$ becomes a random
surface, realised more often near high-density regions of the
native measure.  Consequently, $S(\theta,\omega)$ is a random vector,
and $s_1(\theta,\omega)$, $s_2(\theta,\omega)$ are two random
surfaces.

\begin{figure}[H]
  \centering
  \includegraphics[width=0.85\textwidth]{fig_sv4.png}
  \caption{The Probabilist's sketch (10 June 2026): the two component surfaces
  $s_1(\theta_1,\theta_2)$ (left) and $s_2(\theta_1,\theta_2)$
  (right) of the score vector.  Each surface is bowl-shaped; the
  minimum of each component (green dot) lies at the MLE
  $(\theta_1^*,\theta_2^*)$.  When the data point $\omega$ is random,
  these surfaces are random.}
\end{figure}

\subsection{LM and score test: the one-dimensional case}

We follow notes of June 9, 2026.  We work through the simplest
possible case --- a scalar parameter $\theta\in\mathbb{R}$, simple
null $H_0:\theta=\theta_0$ --- in order to see concretely how the
Lagrangian, the score, and the LM statistic fit together, even when
the problem is trivially solvable without any of this machinery.

\medskip
\noindent\textbf{The log-likelihood as a parabola.}
For $X_1,\ldots,X_n\overset{\mathrm{iid}}{\sim}\mathcal{N}(\theta,1)$
the log-likelihood (up to a constant) is
\[
  \ell(\theta) = -\tfrac{n}{2}(\theta-\bar x)^2,
\]
a downward-opening parabola in $\theta$ with vertex at
$\hat\theta_{\mathrm{MLE}}=\bar x$.  No constraints --- no Lagrangian
needed for unconstrained maximisation.

\begin{figure}[H]
  \centering
  \includegraphics[width=0.55\textwidth]{fig_lms1.png}
  \caption{Sketch: the log-likelihood $\ell(\theta)$ (red parabola)
  as a function of $\theta$.  The vertex is at
  $\hat\theta_{\mathrm{MLE}}=\bar x$.  The blue curve is an
  alternative sample realisation.}
\end{figure}

\medskip
\noindent\textbf{Adding the null hypothesis $\theta=\theta_0$.}
We mark $\theta_0$ on the same graph and read off $\ell(\theta_0)$
(green dot, value $-a$ for some $a\geq 0$).  The constrained
``argmax'' under $H_0:\theta=\theta_0$ is trivially $\tilde\theta=\theta_0$
--- the constraint $\{g=0\}=\{\theta_0\}$ is a single point, so no
optimisation is needed.

\begin{figure}[H]
  \centering
  \includegraphics[width=0.60\textwidth]{fig_lms2.png}
  \caption{Sketch: $\theta_0$ marked on the $\theta$-axis; the green
  dot shows $\ell(\theta_0)=-a$.  The constrained maximum is at this
  single point.}
\end{figure}

\medskip
\noindent\textbf{The constraint as a subset of $\mathbb{R}$.}
In the scalar case the constraint $\{g(\theta)=0\}$ is generically a
finite set of points on the real line.  For
$g(\theta)=\theta-\theta_0$ it is the single point $\theta_0$.

\begin{figure}[H]
  \centering
  \includegraphics[width=0.70\textwidth]{fig_lms3.png}
  \caption{Sketch: the real line with $\theta_0$ marked.  The
  feasible set $\{g=0\}$ is this single point; the constraint reduces
  the one-dimensional domain to a zero-dimensional one.}
\end{figure}

\medskip
\noindent\textbf{The constraint function $g$ and its negation.}
The function $g(\theta)=\theta-\theta_0$ is an affine function of
$\theta$ (left panel below).  Its negation $-g(\theta)=\theta_0-\theta$
is equally valid as a constraint function: $\{g=0\}=\{-g=0\}$.
However, the two choices give \emph{different} Lagrangians
$\ell - \lambda g$ and $\ell - \lambda(-g)$, related by
$\lambda_1 = -\lambda_2$.  The zero set is the same; only the sign
convention for $\lambda^*$ changes.

\begin{figure}[H]
  \centering
  \includegraphics[width=0.75\textwidth]{fig_lms4.png}
  \caption{Left: $g(\theta)=\theta-\theta_0$ (red line, zero at
  $\theta_0$).  Right: $-g(\theta)=\theta_0-\theta$.  Both have the
  same zero set $\{\theta_0\}$; they produce Lagrangians whose
  multipliers differ in sign.}
\end{figure}

\medskip
\noindent\textbf{The Lagrangian in 3d.}
Even in the scalar case we can write the Lagrangian
$\mathcal{L}(\theta,\lambda)=\ell(\theta)-\lambda\,g(\theta)$, which
is a function of two variables $(\theta,\lambda)\in\mathbb{R}^2$.
The term $-\lambda g(\theta)=-\lambda(\theta-\theta_0)$ is a
\emph{bilinear} (saddle-shaped) surface $z=xy$ (after recentring).
Its rulings --- the generating straight lines --- are the families
\[
  \{z = x\,y_0,\ x\in\mathbb{R}\}\quad\text{and}\quad
  \{z = x_0\,y,\ y\in\mathbb{R}\}
\]
for fixed $x_0,y_0\in\mathbb{R}$.  One such ruling is shown in blue
on the red surface below.

\begin{figure}[H]
  \centering
  \includegraphics[width=0.50\textwidth]{fig_lms5.png}
  \caption{The surface $z=\lambda\cdot g(\theta)=\lambda(\theta-\theta_0)$
  (red saddle surface); one ruling (generating line) shown in blue.
  This is the ``LM component'' of the Lagrangian.}
\end{figure}

\medskip
\noindent\textbf{The full Lagrangian in 3d.}
Adding $\ell(\theta)$ (a univariate parabola, constant in $\lambda$,
so it appears as a ``cylindrical'' ridge in 3d) to the saddle surface
$-\lambda g(\theta)$ gives the full Lagrangian surface
$\mathcal{L}(\theta,\lambda)$.  The red parabolic ridge is
$\ell(\theta)$; the blue plane is the LM component at a specific
$\lambda$; the green segment marks the stationary point
$(\tilde\theta,\lambda^*)=(\theta_0,\ell'(\theta_0))$.

\begin{figure}[H]
  \centering
  \includegraphics[width=0.72\textwidth]{fig_lms6.png}
  \caption{The Lagrangian $\mathcal{L}(\theta,\lambda)$ in 3d.  Red:
  the parabolic log-likelihood ridge $\ell(\theta)$, constant in
  $\lambda$.  Blue: the LM component $-\lambda g(\theta)$ (tilted
  plane).  Green: the stationary point at $\tilde\theta=\theta_0$,
  $\lambda^*=\ell'(\theta_0)=n(\bar x-\theta_0)$.  If $H_0$ is true,
  $\lambda^*\approx 0$ and the Lagrangian is nearly flat in $\theta$
  at $\theta_0$.}
\end{figure}

\noindent\textbf{Summary.}
In the scalar one-parameter case the constrained argmax is trivial
(a single point), and the Lagrangian is not needed for solving the
problem.  Nevertheless the formalism applies: the stationary-point
condition $\partial_\theta\mathcal{L}=0$ gives
$\lambda^*=\ell'(\theta_0)=s(\theta_0)$, the score at the null.
This is the LM/score test statistic (before normalisation by
Fisher information).  The LM test and the score test are the same
statistic, and in this univariate case they coincide with the Wald
test as well.

\bigskip

%--------------------------------------------------------------------
\subsection{Full visualisation: normal example with explicit constants}
\label{subsec:lm_normal_full}
%--------------------------------------------------------------------

We now draw everything with explicit numbers from the normal example.
Take $x_1 = 1$ (single observation), true parameter $\theta_0 = 2$,
and the constraint $g(\theta) = \theta - \theta_0 = \theta - 2$.
The log-likelihood is
\[
  \ell(\theta) = -\tfrac{1}{2}\ln(2\pi) - \tfrac{1}{2}(1-\theta)^2,
\]
and the Lagrangian is
\[
  \mathcal{L}(\theta,\lambda)
  = -\tfrac{1}{2}\ln(2\pi) - \tfrac{1}{2}(1-\theta)^2 - \lambda(\theta - 2).
\]

\paragraph{The full Lagrangian and the stationary point.}
Setting $\nabla\mathcal{L} = 0$:
\[
  \begin{cases}
    \dfrac{\partial\mathcal{L}}{\partial\theta}
    = -\tfrac{1}{2}\cdot 2(1-\theta)(-1) - \lambda
    = (1-\theta) - \lambda = 0,\\[6pt]
    \dfrac{\partial\mathcal{L}}{\partial\lambda}
    = -(\theta - 2) = 0.
  \end{cases}
\]
The second equation gives $\theta^* = 2$; substituting into the first:
$1 - 2 - \lambda = 0$, so $\lambda^* = -1$.
The stationary point is $(\theta^*, \lambda^*) = (2, -1)$.

Reading off: $\lambda^* = \ell'(\theta_0) = (1 - \theta_0) = 1 - 2 = -1$
--- the score at the null $\theta_0 = 2$ evaluated at the single
observation $x_1 = 1$.  This is negative because $x_1 < \theta_0$:
the data point lies below the null, so the log-likelihood slopes
downward at $\theta_0$.

\begin{figure}[H]
  \centering
  \includegraphics[width=0.72\textwidth]{fig_lms_full_lagrangian_saddle.png}
  \caption{The full Lagrangian $\mathcal{L}(\theta,\lambda)
  = -\tfrac{1}{2}\ln(2\pi) - \tfrac{1}{2}(1-\theta)^2
  - \lambda(\theta-2)$ for the normal example ($x_1=1$, $\theta_0=2$).
  The blue marker indicates the stationary point at
  $(\theta^*,\lambda^*)=(2,-1)$.
  The surface looks very similar to the $\lambda g$ term alone:
  $\lambda g$ sets the dominant saddle shape, and adding the
  univariate $\ell(\theta)$ (a parabolic ridge, constant in $\lambda$)
  merely shifts the saddle point from $(2,0)$ to $(2,-1)$.
  The value $\lambda^*=-1 = \ell'(\theta_0) = x_1 - \theta_0$
  is the score at the null --- the LM test statistic before
  normalisation by Fisher information.}
\end{figure}

\begin{remark}
The observation that the full Lagrangian looks very similar to the
$\lambda g$ component is not a coincidence for this example.
$\lambda g(\theta)=\lambda(\theta-2)$ is bilinear in $(\theta,\lambda)$
and dominates the shape; the correction $\ell(\theta)$ is a function
of $\theta$ alone (a parabola, independent of $\lambda$), so it enters
as a $\lambda$-invariant ridge that bends the $\theta$-sections without
affecting the saddle structure.  In general, when $\ell$ is approximately
quadratic near $\theta_0$, this additive decomposition
$\mathcal{L} = \ell + (-\lambda g)$ cleanly separates shape (saddle, from
$\lambda g$) and location shift (from $\ell$).
\end{remark}

%--------------------------------------------------------------------
\subsection{The Lagrange multiplier at the saddle equals the score at the null}
\label{subsec:lm_lemma}
%--------------------------------------------------------------------

The two worked examples (normal and exponential) both gave
$\lambda^* = \ell'(\theta_0)$.  This is not a coincidence —
it is a direct consequence of the stationarity conditions for
the Lagrangian, and holds for any differentiable log-likelihood
with any simple null hypothesis.

\paragraph{The identity.}
Take any differentiable function $f:\mathbb{R}\to\mathbb{R}$,
a point constraint $x = x_0$, and the Lagrangian
\[
  \mathcal{L}(x,\lambda) = f(x) - \lambda\,(x - x_0).
\]
At any stationary point $\nabla\mathcal{L} = 0$:
\[
  \frac{\partial\mathcal{L}}{\partial\lambda} = -(x - x_0) = 0
  \implies x^* = x_0,
  \qquad
  \frac{\partial\mathcal{L}}{\partial x} = f'(x) - \lambda = 0
  \implies \lambda^* = f'(x_0).
\]
So the Lagrange multiplier at the saddle equals the derivative of $f$
at the constraint point — period.

Applied to the log-likelihood $\ell(\theta)$ with null $H_0:\theta=\theta_0$:
\[
  \boxed{\lambda^* = \ell'(\theta_0) = s(\theta_0).}
\]
The score at the null \emph{is} the Lagrange multiplier at the
saddle point of the Lagrangian.  The LM test checks whether this
multiplier is close to zero: a large $|\lambda^*|$ means
the log-likelihood is steeply sloped at $\theta_0$, so the data
pull strongly away from the null.

\begin{remark}
Existence: the stationary point with $x^* = x_0$ always exists
as long as $f'(x_0)$ is defined.  Uniqueness of $\lambda^*$
follows immediately: $\lambda^* = f'(x_0)$ is a single number.
If $\ell(\theta)$ is monotone (no interior MLE), there is still a
well-defined stationary point $(\theta_0,\, \ell'(\theta_0))$ of the
Lagrangian, and the LM test still makes sense — it simply registers
a large slope at $\theta_0$.
\end{remark}

\paragraph{Graphical illustration.}
The identity $\lambda^* = s(\theta_0)$ is hard to illustrate
directly because it equates a \emph{slope} on the 2-d
$\ell(\theta)$ graph with a \emph{coordinate} on the 3-d
Lagrangian surface.  The two figures below show both views
simultaneously for the exponential example ($x_1=2$, $\theta_0=1$,
$\lambda^* = s(\theta_0) = 1/\theta_0 - x_1 = -1$).

\begin{figure}[H]
  \centering
  \includegraphics[width=0.65\textwidth]{fig_score_slope_2d.png}
  \caption{2-d view: the log-likelihood $\ell(\theta) = \ln\theta - 2\theta$
  for $x_1 = 2$.  The black tangent line at $\theta_0 = 1$ has slope
  $\ell'(1) = 1/1 - 2 = -1$.  The red right-triangle shows rise$/$run
  $= -1$: one unit right, one unit down.  This slope is $\lambda^*$.}
\end{figure}

\begin{figure}[H]
  \centering
  \includegraphics[width=0.82\textwidth]{fig_lambda_saddle_3d.png}
  \caption{3-d view: the full Lagrangian surface
  $\mathcal{L}(\theta,\lambda) = \ln\theta - 2\theta - \lambda(\theta-1)$.
  The red dot marks the point on the surface directly above the
  saddle; the blue dot on the floor is the saddle projection
  $(\theta^*, \lambda^*) = (1, -1)$, connected by a blue dashed
  vertical.  The red horizontal line at $\lambda = -1$ on the
  $\lambda$-axis identifies the multiplier value.
  Reading off the floor: the $\lambda$-coordinate of the saddle
  is $-1$ — the same number as the slope in the 2-d picture.}
\end{figure}

%--------------------------------------------------------------------
\subsection{Why ``LM test'' and ``score test'' are two names for one thing}
\label{subsec:lm_naming}
%--------------------------------------------------------------------

\paragraph{The score test stands on its own.}
The score test does not need Lagrangians to be defined, computed,
or justified.  Given a log-likelihood $\ell(\theta)$ and a null
$H_0:\theta=\theta_0$, the test statistic is simply
\[
  \xi_{LM}
  = \frac{[s(\theta_0)]^2}{\mathcal{I}(\theta_0)}
  = \frac{[\ell'(\theta_0)]^2}{\mathcal{I}(\theta_0)}
  \xrightarrow{d} \chi^2_1 \quad \text{under } H_0.
\]
The score $s(\theta_0)$ is evaluated \emph{at the null} $\theta_0$,
not at the MLE $\hat\theta$ (where $s(\hat\theta)=0$ by definition,
making the statistic trivially zero).
The Fisher information $\mathcal{I}(\theta_0)$ in the denominator
normalises the score to unit variance, so that the squared
standardised score has a $\chi^2_1$ limit by the CLT.

\paragraph{Where the name ``LM'' comes from.}
Silvey (1959) derived the same statistic via a Lagrangian argument:
constrain $\ell(\theta)$ to $H_0$ using a multiplier $\lambda$,
find the saddle point of the Lagrangian, and read off $\lambda^*$.
The result of this section shows that $\lambda^* = s(\theta_0)$,
so the Lagrange multiplier \emph{is} the score at the null.
Engle (1984) proved the asymptotic equivalence of all three tests
(Wald, LR, Score/LM) and established the name ``LM test'' alongside
``score test'' as synonyms.  The Lagrangian machinery adds no
computational content --- it only explains the name.

\paragraph{Shadow price interpretation (a remark on language).}
$\lambda^*$ is sometimes called the \emph{shadow price} of the
constraint.  In the inequality-constraint literature ``relaxation''
means moving from $g\le 0$ to $g\le\varepsilon$.  For our equality
constraint $\theta=\theta_0$ there is no relaxation in that sense;
instead the shadow price measures the rate of change of the
constrained optimum when the constraint level is \emph{shifted}:
\[
  \frac{d}{d\varepsilon}
  \max_{\theta:\,\theta=\theta_0+\varepsilon}\ell(\theta)
  \bigg|_{\varepsilon=0}
  = \ell'(\theta_0) = \lambda^*.
\]
A large $|\lambda^*|$ means the constrained optimum is sensitive
to where exactly we place $\theta_0$: the data push hard away from
the null, and the LM test detects this.

%--------------------------------------------------------------------
\subsection{Worked example: exponential distribution}
\label{subsec:lm_exp}
%--------------------------------------------------------------------

We repeat the full LM story for the exponential family, which has
non-quadratic log-likelihood and provides a useful contrast to the
normal case.

\paragraph{Setup.}
The exponential density with rate parameter $\theta > 0$ is
\[
  f_\theta(x) = \mathbf{1}_{x\ge 0}\cdot\theta\,e^{-\theta x}.
\]

\begin{figure}[H]
  \centering
  \includegraphics[width=0.55\textwidth]{fig_exp1_density.png}
  \caption{Exponential density $f_\theta(x) = \theta e^{-\theta x}$
  for $x\ge 0$.  The density starts at height $\theta$ and decays
  exponentially; the $y$-intercept equals the rate parameter.}
\end{figure}

\paragraph{Likelihood and log-likelihood surfaces.}
For a single observation $x_1$ the likelihood is
$L_x(\theta) = \theta e^{-\theta x}$ and the log-likelihood is
\[
  \ell_x(\theta) = \ln\theta - \theta x.
\]

\begin{figure}[H]
  \centering
  \includegraphics[width=0.72\textwidth]{fig_exp2_lik3d.png}
  \caption{The likelihood surface $z = \theta e^{-\theta x}$
  over the $(\theta, x)$ plane ($x=\theta$ axis, $y=x_1$ axis).
  The surface rises steeply near $\theta=0$ for small $x_1$ and
  decays in both directions.}
\end{figure}

\begin{figure}[H]
  \centering
  \includegraphics[width=0.72\textwidth]{fig_exp3_loglik3d.png}
  \caption{The log-likelihood surface $z = \ln\theta - \theta x$
  over the $(\theta, x)$ plane.  The surface goes to $-\infty$
  as $\theta\to 0^+$ (from the $\ln\theta$ term) and to $-\infty$
  as $\theta\to+\infty$ (from the $-\theta x$ term), so each
  $x_1$-section has a unique interior maximum.  The plot window
  must be restricted to small $\theta$ (e.g.\ $\theta\in(0,3]$)
  to make the $\ln\theta\to-\infty$ singularity visible; at large
  $\theta$ it is compressed into a narrow strip near the left edge.}
\end{figure}

\paragraph{MLE.}
Maximising $\ell_x(\theta) = \ln\theta - \theta x$ over $\theta > 0$:
\[
  \frac{d\ell}{d\theta} = \frac{1}{\theta} - x = 0
  \implies \hat\theta_{\mathrm{MLE}} = \frac{1}{x}.
\]
The MLE is the reciprocal of the observation.  As a function of $x_1$
it is a hyperbola: larger observations shift the rate estimate downward.

\begin{figure}[H]
  \centering
  \includegraphics[width=0.55\textwidth]{fig_exp4_mle2d.png}
  \caption{The MLE $\hat\theta = 1/x_1$ as a function of the
  observation $x_1 > 0$.  The curve is a rectangular hyperbola
  with vertical asymptote at $x_1=0$.}
\end{figure}

\paragraph{Specific case: $x_1 = 2$.}
With $x_1 = 2$ the MLE is $\hat\theta = 1/2$, and the log-likelihood
reduces to
\[
  \ell(\theta)\big|_{x_1=2} = \ln\theta - 2\theta.
\]

\begin{figure}[H]
  \centering
  \includegraphics[width=0.72\textwidth]{fig_exp5_mle3d.png}
  \caption{The log-likelihood surface $z = \ln\theta - \theta x_1$
  with the MLE ridge $\hat\theta = 1/x_1$ marked (blue curve on the
  surface).  Each vertical cross-section at fixed $x_1$ achieves its
  maximum at $\theta = 1/x_1$; the ridge traces the argmax as $x_1$
  varies.}
\end{figure}

\begin{figure}[H]
  \centering
  \includegraphics[width=0.72\textwidth]{fig_exp6_loglik2d.png}
  \caption{The one-dimensional log-likelihood $\ell(\theta)=\ln\theta-2\theta$
  for $x_1=2$.  The function goes to $-\infty$ as $\theta\to 0^+$,
  rises to its maximum at $\hat\theta = 1/2$, then decreases to
  $-\infty$ as $\theta\to+\infty$.  Argmax: $\theta = 1/2$.}
\end{figure}

\paragraph{Lagrangian setup.}
Set the null $\theta_0 = 1$ and the constraint
$g(\theta) = \theta - \theta_0 = \theta - 1$.
The Lagrangian is
\[
  \mathcal{L}(\theta,\lambda)
  = \underbrace{\ln\theta - 2\theta}_{\ell(\theta)}
    - \lambda(\theta - 1).
\]
The second term alone is $T = \lambda(\theta-1) = \lambda\theta - \lambda$,
a bilinear (saddle) surface.

\begin{figure}[H]
  \centering
  \includegraphics[width=0.72\textwidth]{fig_exp7_loglik_lambda3d.png}
  \caption{The log-likelihood surface $z = \ln\theta - 2\theta$
  extended trivially in the $\lambda$ direction (constant in $\lambda$).
  This is the $\ell(\theta)$ component of the Lagrangian viewed
  in $(\theta,\lambda)$ space: a ridge parallel to the $\lambda$-axis
  with peak at $\hat\theta = 1/2$.}
\end{figure}

\paragraph{The constraint term and its saddle.}
The second term $T(\theta,\lambda) = \lambda(\theta-1)$ is a bilinear
surface.  Setting $\nabla T = 0$:
\[
  \begin{cases}
    \dfrac{\partial T}{\partial\theta} = \lambda = 0,\\[4pt]
    \dfrac{\partial T}{\partial\lambda} = \theta - 1 = 0,
  \end{cases}
  \implies (\theta^*, \lambda^*) = (1,\, 0).
\]
The saddle of the constraint term alone sits at $(\theta_0, 0)$; adding
$\ell(\theta)$ will shift the $\lambda$-coordinate of the stationary
point, while leaving $\theta^* = \theta_0$ unchanged.

\begin{figure}[H]
  \centering
  \includegraphics[width=0.72\textwidth]{fig_exp8_constraint3d.png}
  \caption{The constraint term $T(\theta,\lambda) = \lambda(\theta-1)$
  as a 3-d surface (familiar bilinear saddle).
  Saddle point at $(\theta,\lambda) = (1, 0)$ with $T = 0$.}
\end{figure}

\paragraph{Full Lagrangian and stationary point.}
Setting $\nabla\mathcal{L} = 0$:
\[
  \begin{cases}
    \dfrac{\partial\mathcal{L}}{\partial\theta}
      = \dfrac{1}{\theta} - 2 - \lambda = 0,\\[8pt]
    \dfrac{\partial\mathcal{L}}{\partial\lambda}
      = -(\theta - 1) = 0.
  \end{cases}
\]
The second equation gives $\theta^* = 1 = \theta_0$.
Substituting into the first: $1/1 - 2 - \lambda = 0$, so
$\lambda^* = -1$.

The stationary point is $(\theta^*, \lambda^*) = (1,\,-1)$.
Crucially, $\lambda^* = -1 = \ell'(\theta_0) = 1/\theta_0 - x_1
= 1/1 - 2 = -1$: the Lagrange multiplier at the constrained
optimum equals the score evaluated at the null.
This is the LM test statistic (before normalisation).

The Fisher information for $\mathrm{Exp}(\theta)$ is
$\mathcal{I}(\theta) = 1/\theta^2$, so at $\theta_0 = 1$:
$\mathcal{I}(\theta_0) = 1$.  The normalised LM statistic is
\[
  \xi_{LM}
  = \frac{[\ell'(\theta_0)]^2}{\mathcal{I}(\theta_0)}
  = \frac{(-1)^2}{1} = 1.
\]

\begin{figure}[H]
  \centering
  \includegraphics[width=0.72\textwidth]{fig_exp9_lagrangian3d.png}
  \caption{The full Lagrangian $z = \ln\theta - 2\theta - \lambda(\theta-1)$
  for $x_1=2$, $\theta_0=1$.  The surface looks dominated by the
  univariate $\ell(\theta)$ component (the ridge shape), making
  the saddle point hard to read off visually.}
\end{figure}

\begin{figure}[H]
  \centering
  \includegraphics[width=0.72\textwidth]{fig_exp10_lagrangian_annotated.png}
  \caption{The full Lagrangian surface annotated with the stationary
  point $(\theta^*, \lambda^*) = (1, -1)$ (blue marker on the floor).
  The saddle is invisible from visual inspection alone; finding it
  requires solving $\nabla\mathcal{L}=0$ by hand.
  The $\lambda^* = -1$ is the score $\ell'(\theta_0) = 1/\theta_0 - x_1
  = -1$: negative because $x_1 = 2 > 1/\theta_0 = 1$, i.e.\ the
  observation exceeds the reciprocal of the null rate, pushing the
  MLE below $\theta_0$ and putting the log-likelihood on a downward
  slope at $\theta_0$.}
\end{figure}

\begin{remark}
  Comparison with the normal example.  In both cases
  $\lambda^* = \ell'(\theta_0) = s(\theta_0)$, the score at the
  null.  For the normal case ($x_1=1$, $\theta_0=2$) we obtained
  $\lambda^* = x_1 - \theta_0 = -1$.
  For the exponential case ($x_1=2$, $\theta_0=1$) we obtain
  $\lambda^* = 1/\theta_0 - x_1 = -1$.
  The numerical values coincide by coincidence; the structure is
  universal: $\lambda^* = s(\theta_0)$ always, and the LM test
  checks whether this value is close to zero.
\end{remark}

%--------------------------------------------------------------------
\subsection{Two questions on Lagrangian theory}
\label{subsec:lm_theory_questions}
%--------------------------------------------------------------------

Two conceptual questions arose while working through the examples
above.  Both concern what can be said about the Lagrangian's
stationary point \emph{without} doing the computation --- purely
from the shape of $f$ and $g$.

\paragraph{Question 1: directional behaviour at the saddle.}
The Lagrangian $\mathcal{L}(\theta,\lambda)=f(\theta)-\lambda g(\theta)$
has no local optima of its own --- only a stationary point, which is
generically a genuine saddle (one direction up, one direction down).
Can the directional behaviour be read off in advance, from $f$ and
$g$ alone, without solving for $(\theta^*,\lambda^*)$?

\emph{An important correction first.}
In the scalar examples used throughout this section ($\theta\in\mathbb{R}$,
single point constraint $g(\theta)=\theta-\theta_0$), the constraint
set $\{g=0\}=\{\theta_0\}$ is a single point: a $0$-dimensional
manifold.  There is no room to ``move along the constraint,'' so
the restricted function $\tilde f$ (Question 2's object) has no
interior direction to differentiate.  To ask a meaningful question
about \emph{directional} behaviour of $\tilde f$ near the
constrained optimum, we need at least one free direction left on
the constraint set --- i.e.\ a vector parameter
$\theta=(\theta_1,\theta_2)\in\mathbb{R}^2$ with a single scalar
constraint $g(\theta_1,\theta_2)=0$, exactly the setup of
the Probabilist's original Lagrange-multiplier pictures (\S\ref{subsec:basic-lm}--%
\ref{subsec:lm_gradient}: $f(x,y)$, constraint curve $\{g=0\}$).
We work in that setting here, writing $\theta=(x,y)$ to match those
figures.

\emph{Part 1: every Lagrangian is affine in each multiplier.}
For $k$ constraints $g_1,\ldots,g_k$, the Lagrangian
$\mathcal{L}=f-\sum_j \lambda_j g_j$ is affine in each $\lambda_j$
separately (holding $\theta$ and the other multipliers fixed):
$\lambda_j \mapsto f(\theta) - \lambda_j g_j(\theta) - (\text{const.})$.
An affine function's graph is a straight line, so the generating
curves of the Lagrangian surface along any $\lambda_j$-direction
are always straight lines (the ``rulings'' of the saddle seen in
Figure~26).  At a stationary point, $g_j(\theta^*)=0$ for every $j$
(the constraints hold), so each of these affine sections has slope
$-g_j(\theta^*) = 0$: the line is horizontal.  This confirms and
generalises the earlier finding: the flat ``seam'' along each
$\lambda_j$-direction through the solution is not a coincidence
of the one-constraint case --- it holds for any number of
constraints, by the same one-line argument.

\emph{Part 2: fixing $\lambda=\lambda^*$ and varying $\theta$ near $\theta^*$.}
Now fix $\lambda=\lambda^*$ and ask about the $\theta$-behaviour of
$\mathcal{L}(\cdot,\lambda^*)$ near $\theta^*$, in relation to the
restricted function $\tilde f$ on the constraint curve
$\{g=0\}$ introduced in Question~2.  Write $\tilde\theta^*$ for the
point on the curve corresponding to $\theta^*$ (so
$\tilde f(\tilde\theta^*) = f(\theta^*)$).  Since $g(\theta^*)=0$,
\[
  f(\theta^*) = \tilde f(\tilde\theta^*) = \mathcal{L}(\theta^*,\lambda^*).
\]
The three quantities agree \emph{at the point}, but this does not
mean $f$, $\tilde f$, and $\mathcal{L}(\cdot,\lambda^*)$ behave alike
\emph{around} the point --- they are functions on different domains
($f$ and $\mathcal{L}(\cdot,\lambda^*)$ on the full ambient space,
$\tilde f$ only on the curve), and conflating them is an error worth
flagging explicitly, since it is easy to make.

First derivatives.  By construction (the stationarity condition),
\[
  \nabla_\theta\mathcal{L}(\theta^*,\lambda^*)
  = \nabla f(\theta^*) - \lambda^*\nabla g(\theta^*) = 0.
\]
This says $\nabla f(\theta^*) = \lambda^*\nabla g(\theta^*)$ ---
\emph{not} that $\nabla f(\theta^*)=0$.  In general $f$ itself has a
nonzero gradient at $\theta^*$ (it is only the projection along the
constraint direction that vanishes).  So the unconstrained $f$ can
have any slope at $\theta^*$; nothing here forces $f$ to be flat
there.

What about $\tilde f\,'$, the intrinsic derivative of $f$ restricted
to the curve $\{g=0\}$?  By definition of constrained stationarity
(Question~2), $\theta^*$ is a stationary point of $\tilde f$ exactly
when $\nabla f(\theta^*)\parallel\nabla g(\theta^*)$ --- which is
precisely the condition we just derived from $\nabla_\theta
\mathcal{L}=0$.  So $\tilde f\,'(\tilde\theta^*)=0$ always, at any
stationary point of the Lagrangian: this is one honest fact we can
state about $\tilde f$ near $\theta^*$, in contrast to $f$ itself
(which need not be flat there).

Second derivatives are where the behaviours can genuinely differ.
$\tilde f$ lives on a curve, so $\tilde f\,''(\tilde\theta^*)$ is a
single number determining whether $\tilde\theta^*$ is a constrained
local max, min, or inflection of $\tilde f$ --- and this number
depends on the curvature of \emph{both} $f$ and $g$ along the curve
(the bordered Hessian), not on $f''$ alone.  Meanwhile
$\mathcal{L}(\cdot,\lambda^*)$, as a function on the full
2-dimensional $\theta$-space, has its own Hessian
$\nabla^2 f(\theta^*) - \lambda^*\nabla^2 g(\theta^*)$, a $2\times 2$
matrix --- a strictly richer object than the scalar
$\tilde f\,''(\tilde\theta^*)$.  The relation between them is exactly
the content of the bordered-Hessian test: $\tilde\theta^*$ is a
constrained local max of $\tilde f$ iff the bordered Hessian
\[
  \bar H = \begin{pmatrix} 0 & \nabla g(\theta^*)^\top \\
  \nabla g(\theta^*) & \nabla^2 f(\theta^*)-\lambda^*\nabla^2 g(\theta^*)
  \end{pmatrix}
\]
has the appropriate sign pattern on the subspace orthogonal to
$\nabla g(\theta^*)$ --- a strictly weaker requirement than
$\nabla^2 f(\theta^*)-\lambda^*\nabla^2g(\theta^*)$ being negative
definite on the whole space.  So: the second-order behaviour of
$\tilde f$ cannot in general be read off from $f''$ alone (our
earlier scalar-case shortcut, valid only because there the curve
degenerated to a point and the question of an interior second
derivative did not arise); in the genuine multivariate case it
requires this restricted (bordered) computation.

\emph{Why $\mathcal{L}$ is always a genuine saddle in $(\theta,\lambda)$.}
This part of the original answer survives unchanged, since it
concerned the full Lagrangian $\mathcal{L}(\theta,\lambda)$ jointly
in $(\theta,\lambda)$, not the restricted $\tilde f$.  In the
one-constraint, scalar-multiplier case, the mixed Hessian block is
$-\nabla g(\theta^*)\ne 0$ (constraint qualification), forcing one
positive and one negative eigenvalue overall: $\mathcal{L}$ is a
saddle in the joint $(\theta,\lambda)$ sense regardless of what
$\tilde f\,''$ turns out to be.

\paragraph{Question 2: stationarity on the constrained manifold.}
The unconstrained stationary points of $f$ are irrelevant to the
constrained problem --- but the constrained problem is itself a
function restricted to a lower-dimensional set (a curve, in the
two-variable case), and on that set one can still speak of
\emph{local optima} directly.  Can one also speak of stationarity
intrinsically, without invoking the full-dimensional derivative
that requires a full neighbourhood?

The answer is yes, and it is not a vague analogy: when $g$ is
smooth and $\nabla g\ne 0$ at the point in question (the constraint
qualification), the implicit function theorem guarantees that
$\{g=0\}$ is locally a smooth curve (a 1-dimensional manifold,
in the two-variable case), parametrisable by a single coordinate.
Restricting $f$ to this curve gives an honest one-variable function
$\tilde f$ of that coordinate, with an honest derivative
$\tilde f\,'$ --- not merely something ``directional-derivative-like.''
The stationary points of $\tilde f$ (where $\tilde f\,'=0$) are
precisely the local optima of $f$ on the constraint curve.

The classical Lagrange theorem is exactly the statement that these
intrinsic stationary points of $\tilde f$ coincide with the
projections of the stationary points of $\mathcal{L}$:
\[
  \tilde f\,'(\theta^*) = 0
  \iff
  \nabla f(\theta^*) \parallel \nabla g(\theta^*)
  \iff
  \exists\,\lambda^*:\ \nabla_{\theta,\lambda}\mathcal{L}(\theta^*,\lambda^*)=0.
\]
This is a genuine \emph{if and only if}, not merely an inclusion of
one set of points by the other.  So the ``stronger statement''
anticipated in the question is, in fact, already the standard
Lagrange theorem --- once the intrinsic derivative on the
constraint curve is defined via the implicit function theorem,
there is no missing or ill-defined notion left over: the apparent
gap (``I don't have a good definition'') is closed by the
constraint qualification, which is precisely the hypothesis that
makes $\{g=0\}$ a smooth curve with a well-defined tangent direction
at every point.

%--------------------------------------------------------------------

\end{document}
